
\begin{exercise}[Products of random variables \Coffeecup]
    This entire chapter is about sums of random variables. In this exercise we want
    to demonstrate that these results can be easily translated to products of
    random variables with the help of the logarithm. For
    \((X_n)_{n\in \nat}\) independent and identically distributed random variables
    with \(X_n > 0\) and \(\E{|\ln(X_1)|} < \infty\) prove that the
    `geometric mean' converges almost surely to the `geometric expectation',
    that is'
    \[
        \Bigl(\prod_{i=1}^n X_i \Bigr)^\frac1n
        \overset{\as}\to \exp\bigl(\E[\big]{\ln(X_1)}\bigr).
    \]
\end{exercise}

\begin{solution}
    Observe that the logarithm of the geometric mean converges by the law of
    large numbers to the expectation of the logarithm, that is
    \[
        \ln\Bigl(\Bigl(\prod_{i=1}^n X_i \Bigr)^\frac1n\Bigr)
        = \frac1n \sum_{i=1}^n \ln(X_i)
        \overset{\as}\to \E[\big]{\ln(X_1)}.
    \]
    Since \(\exp(x)\) is continuous the claim follows from continuous mapping.
\end{solution}
